% §16.5 Where the Framework Still Helps
%
% Cross-reference resolution status:
%
% DEFINED here:
%   \label{sec:where-framework-helps}         -- this section
%
% ALREADY DEFINED (chapter 4):
%   \Cref{ax:element-indep}                   -- chapter 4 §4.2 Axiom 1
%   \Cref{ax:block-rates-indep}               -- chapter 4 §4.3 Axiom 2
%
% ALREADY DEFINED (chapter 6):
%   \Cref{sec:perfect-filter-line}            -- chapter 6 §6.4
%   \Cref{def:perfect-filter-line}            -- chapter 6 §6.4
%
% ALREADY DEFINED (chapter 9):
%   \Cref{sec:n-ary-channels}                 -- chapter 9 §9.2
%   \Cref{def:n-ary-channel}                  -- chapter 9 §9.2
%   \Cref{sec:n-ary-kronecker}                -- chapter 9 §9.3
%   \Cref{thm:n-ary-kronecker}                -- chapter 9 §9.3
%
% ALREADY DEFINED (chapter 7):
%   \Cref{def:ppv}                            -- chapter 7 §7.4 PPV definition
%
% ALREADY DEFINED (chapter 11):
%   \Cref{thm:chernoff}                       -- chapter 11 §11.4 Chernoff bound
%
% INTERNAL refs (chapter 16):
%   \Cref{sec:varying-rates}                  -- §16.2
%   \Cref{sec:correlated-noise}               -- §16.3
%   \Cref{sec:adversarial-inputs}             -- §16.4
%   \Cref{prop:azuma-hoeffding}               -- §16.2
%   \Cref{prop:mcdiarmid}                     -- §16.3
%
% §16.5 pushes back against pure relinquishment: even when Axiom 1
% or Axiom 2 fails, the framework's geometric and channel-matrix
% intuitions still apply, often with minor modifications. Catalogue
% the perfect-filter line, the channel matrix as a sequence, the PPV
% per-query formula, and close with a pragmatic rule about replacing
% iid concentration tools with their non-iid counterparts.

\section{Where the Framework Still Helps}
\label{sec:where-framework-helps}

The three preceding sections enumerated systems for which strict
iid analysis breaks. A reasonable reading at this point might be
that the Bernoulli framework, having been built atop two axioms,
is useless once either axiom fails. The reading is too pessimistic.
The framework's geometric and structural intuitions reach beyond
their strict iid statement, and several of its constructs survive
the failures of Axiom 1 or Axiom 2 with only minor accommodation.
This section catalogues the partial successes.

\subsection*{Geometric reach: the perfect-filter line}

\Cref{sec:perfect-filter-line} introduced the perfect-filter line
as the set of systems in $[0, 1]^2$ rate-plane that achieve either
$\fprate = 0$ or $\fnrate = 0$ exactly, identified by their
position on one of the two axes. The classification was set up for
systems with stable rates; \Cref{def:perfect-filter-line} reads
the perfect-filter property off a single channel matrix.

The classification survives the failures of \Cref{sec:varying-rates}
intact. A system whose rates drift over time still has a
$(\fprate(t), \fnrate(t))$ point at each query, and the trajectory
in $[0, 1]^2$ traces a curve. The perfect-filter line still
classifies each individual point: a system whose
$(\fprate(t), \fnrate(t))$ trajectory stays on the $\fprate = 0$
axis throughout is one-sided in the same sense as before, even
though its $\fnrate(t)$ is not constant. Bloom filters past
capacity, for example, remain $\fprate = 0$ never (they have
nonzero false-positive rate at every time) but stay on the
$\fnrate = 0$ axis at every time (a Bloom filter never produces a
false negative).

The classification also survives \Cref{sec:correlated-noise}. The
joint distribution between queries shifts away from the iid
product, but the marginal at each query has a well-defined
$(\fprate, \fnrate)$ point, and the perfect-filter classification
applies per-query. A correlation between two filter queries does
not move either filter's marginal rates off whichever axis they
were on; it only changes how the joint behaves around those
marginals.

The geometric statement, then, is portable: the perfect-filter
line classifies systems by per-query marginals, and per-query
marginals are well-defined even when their iid-product
extrapolation is not. A system designer can use the perfect-filter
classification to reason about whether a noisy bit ever errs
toward $\top$, ever errs toward $\bot$, or both, even in the face
of drift or correlation. What disappears is the right to chain
per-query rates into multi-query bounds without further work; the
classification itself remains.

\subsection*{Channel-matrix reach: $Q$ becomes a sequence $(Q_t)$}

\Cref{sec:n-ary-channels} introduced the channel matrix $Q$ as
the per-block error matrix; \Cref{def:n-ary-channel} fixed its
entries as the conditional probabilities of observed labels given
true labels. Under Axiom 2, $Q$ is a single matrix; under
\Cref{sec:varying-rates}'s failure, $Q$ becomes a sequence
$(Q_t)_{t = 1, \ldots, m}$, with one matrix per query.

The channel-matrix framework absorbs this generalization with
little friction. A sequence of matrices $(Q_t)$ still describes
the per-query input-output behavior; what fails is the stability
claim that lets one treat a sequence of queries as draws from a
single matrix. Many derived properties survive at the per-query
level. The eigenstructure of $Q_t$, the channel capacity at query
$t$, and the Bayes-optimal decision rule given the matrix and a
prior all read off $Q_t$ alone.

\Cref{sec:n-ary-kronecker}'s Kronecker theorem
(\Cref{thm:n-ary-kronecker}) factored independent channels into
Kronecker products. The theorem also generalizes. When the
component channels have time-varying or correlated rates, the
Kronecker factorization applies conditionally on the structure of
the dependence: for time-varying rates, factor channel matrices
into Kronecker components at each time step, giving a sequence of
Kronecker products $(Q_t^{(1)} \otimes \cdots \otimes Q_t^{(d)})$
rather than a single product; for correlated channels, the
factorization fails marginally but holds conditional on whatever
shared randomness drives the correlation. The Kronecker geometry,
in other words, applies sliced or conditioned rather than
unconditionally.

\subsection*{PPV/NPV reach: the per-query formula stands}

\Cref{def:ppv} expressed PPV as a function of $\fprate$, $\fnrate$,
and prevalence $\pi$. The formula is per-query: it took a single
$(\fprate, \fnrate, \pi)$ triple and returned the conditional
probability of a true positive given a positive report. Nothing in
the derivation depended on iid sampling across queries; the
identity is structural.

Under any of the failure modes of \Cref{sec:varying-rates},
\Cref{sec:correlated-noise}, or \Cref{sec:adversarial-inputs}, the
PPV formula still holds at each individual query. What changes is
that the triple $(\fprate, \fnrate, \pi)$ may itself be
time-varying or query-dependent. Time-averaging PPV across queries
requires care: the time-average PPV is not generally equal to the
PPV evaluated at the time-average rates, because PPV is nonlinear
in $\pi$. But at any single query, the formula gives the correct
conditional probability given that query's rates and prevalence.

The practical implication is that PPV and NPV remain useful
diagnostic quantities even outside strict iid, provided one
reports them per-query (or per-window) rather than as a single
global summary. A system designer who notices drift in $\pi$ can
plot PPV as a function of time and reason about its trajectory,
without needing to commit to a fixed $\pi$ that the data would
contradict.

\subsection*{A pragmatic rule}

The catalogue above suggests a working rule for analysts faced
with non-iid Bernoulli data. The rule has two halves.

First: keep the framework's geometry and channel-matrix view, even
when Axioms 1 and 2 fail. The perfect-filter line classifies
per-query, the channel matrix factorizes per-query (or per-block
of independent channels), and the PPV/NPV formulas evaluate
per-query. None of these per-query statements assumes iid; all
remain useful diagnostic tools.

Second: replace the iid analytical tools with their non-iid
counterparts. Where \Cref{thm:chernoff} bounded sums of iid
Bernoulli random variables, use \Cref{prop:azuma-hoeffding} for
the martingale-difference setting of \Cref{sec:varying-rates} or
\Cref{prop:mcdiarmid} for the bounded-differences setting of
\Cref{sec:correlated-noise}. Where the binomial confidence
interval gave a sampling-error bound under iid, use the
bounded-difference concentration or the conditional bounds for
the dependent case. Where neither iid nor a tractable dependence
structure applies (\Cref{sec:adversarial-inputs}), drop the
expectation and use worst-case bounds, then close the relevant
side channels via cryptographic or oblivious primitives.

The framework's affirmative content is per-query: a noisy bit
with rates $(\fprate, \fnrate)$, a channel matrix $Q$, a derived
quantity such as PPV computed from those rates. Strict iid
extends this per-query content to multi-query bounds. When iid
fails, the per-query content survives; only the cross-query
machinery requires replacement. Treating the framework as a
collection of per-query objects glued together with iid wrappers,
rather than as an iid edifice, restores most of its utility in
the non-iid setting. The wrappers change; the underlying
geometry does not.
